% Del dual-issue al dual-issue con predicción de saltos (documento independiente).
% Compilar con ./build.sh (genera primero las figuras de ./figs).
\documentclass[11pt,a4paper]{article}
\usepackage[utf8]{inputenc}
\usepackage[T1]{fontenc}
\usepackage[spanish,es-noquoting,es-tabla,es-nolists]{babel}
\usepackage{lmodern}            % typewriter (instrucciones)
\usepackage{mathpazo}           % texto y matemáticas: Palatino
\usepackage[scale=0.9]{tgheros} % sans serif de títulos, pies de figura y figuras
\linespread{1.06}
\usepackage{microtype}
\usepackage[a4paper,margin=2.4cm]{geometry}
\usepackage{graphicx}
\usepackage{pdflscape}
\usepackage{booktabs}
\usepackage{array}
\usepackage{amsmath}
\usepackage{caption}
\usepackage{enumitem}
\usepackage{xcolor}
\usepackage{listings}
\usepackage[most]{tcolorbox}
\usepackage{titlesec}
\usepackage[hidelinks]{hyperref}

\graphicspath{{figs/}}
\definecolor{hhblue}{RGB}{11,111,184}
\definecolor{hzd}{RGB}{217,95,2}
\definecolor{new}{RGB}{21,128,61}
\definecolor{opt}{RGB}{139,90,43}

\captionsetup{labelfont={bf,sf,color=hhblue},textfont={bf,sf,small},labelsep=quad,
  figurename=Figura,tablename=Tabla}

% headings as in the book: sans bold, number in black, title in blue
\titleformat{\section}{\sffamily\bfseries\Large}{\thesection}{0.7em}{\color{hhblue}}
\titleformat{\subsection}{\sffamily\bfseries\large}{\thesubsection}{0.7em}{\color{hhblue}}
\titleformat{\paragraph}[runin]{\sffamily\bfseries\normalsize}{}{0pt}{\color{hhblue}}
\titlespacing*{\section}{0pt}{3.2ex plus 1ex minus .2ex}{1.6ex plus .2ex}
\titlespacing*{\subsection}{0pt}{2.6ex plus 1ex minus .2ex}{1.2ex plus .2ex}
\titlespacing*{\paragraph}{0pt}{1.6ex plus .5ex minus .2ex}{0.8em}

% signal names in italics (as in the book); lane number as subscript
\newcommand{\sig}[1]{\textit{#1}}
\newcommand{\sigk}[2]{\textit{#1}\textsubscript{#2}}
\newcommand{\ins}[1]{\texttt{#1}}

% HDL examples, numbered per chapter
\lstdefinestyle{vlog}{language=Verilog, basicstyle=\ttfamily\footnotesize, columns=fullflexible,
  keywordstyle=\color{hhblue}\bfseries, commentstyle=\color{black!55}\itshape,
  showstringspaces=false, keepspaces=true, upquote=true,
  breaklines=true, breakindent=0pt, breakautoindent=true}
\newcounter{hdl}
\newtcblisting{hdl}[1]{listing only, breakable, colback=black!3, colframe=black!35,
  boxrule=0.4pt, arc=1.5pt, left=4pt, right=4pt, top=2pt, bottom=2pt,
  title={\sffamily\footnotesize\refstepcounter{hdl}\textbf{\textcolor{hhblue}{Ejemplo HDL \thehdl}}\quad #1},
  coltitle=black, colbacktitle=black!10, listing options={style=vlog}}

\setlist{itemsep=1pt, topsep=3pt}
\setlist[itemize]{label=\textbullet}

\begin{document}

\begin{center}
  {\sffamily\bfseries\huge\color{hhblue} Del dual-issue a la predicción de saltos\par}
  \vspace{0.6ex}
  {\sffamily\large Branch Target Buffer con contador de 2 bits en un RISC-V de dos vías\par}
  \vspace{1.2ex}
  \rule{\linewidth}{0.4pt}
\end{center}
\vspace{1ex}

Nuestro procesador dual-issue predice siempre que los saltos no se toman: mientras un salto
avanza hasta Execute, donde se resuelve, Fetch sigue buscando las instrucciones que vienen a
continuación en memoria. Cuando el salto sí se toma, los dos pares que se buscaron detrás
de él se descartan: el salto cuesta dos ciclos, en los que se dejan de emitir hasta cuatro
instrucciones. El salto que cierra un bucle se toma casi en cada vuelta, así que ese costo
se paga casi siempre.

En esta guía reemplazamos esa predicción fija por una \emph{predicción de saltos}
(\emph{branch prediction}) dinámica: en Fetch, antes de saber si el salto se toma, el
procesador predice su resultado según lo que hizo ese mismo salto las veces anteriores, y
sigue buscando instrucciones por el camino predicho. El predictor es un \emph{Branch Target
Buffer} (BTB): una tabla que guarda, para cada salto tomado recientemente, su dirección de
destino y un contador de 2 bits.

La sección~\ref{sec:base} repasa el procesador de partida y muestra cuánto cuesta un salto
tomado. La sección~\ref{sec:overview} presenta la idea de la predicción. A partir de la
sección~\ref{sec:counter}, cada sección agrega una parte al hardware: el contador, la tabla,
la etapa Fetch, el transporte de la predicción, su verificación y los cambios en la Hazard
Unit, hasta llegar al procesador completo de la Fig.~\ref{fig:proc-full}. Algunos bloques se
detallan con fragmentos de Verilog (los \emph{Ejemplos HDL}). Los diagramas de ejecución
provienen de simular ese mismo hardware.

% =====================================================================
\section{El punto de partida: el procesador dual-issue}
\label{sec:base}

Partimos de un RISC-V
superescalar de dos vías y en orden, construido a partir del pipeline de cinco etapas
de~\cite{harris}. En cada ciclo busca y decodifica dos instrucciones
consecutivas, y puede emitir dos, una o ninguna de ellas a Execute. Usamos los mismos
términos que en su diseño:
\begin{itemize}
  \item \textbf{Par}: las dos instrucciones consecutivas que se buscan y decodifican juntas.
  \item \textbf{Slot 0 / slot 1}: posición dentro del par. El slot~0 contiene la
        instrucción \emph{más antigua} y el slot~1 la \emph{más joven}, la que sigue en
        memoria.
  \item \textbf{Vía} (\emph{lane}): cada una de las dos copias del camino de ejecución. El
        slot~0 entra a la vía~0 y el slot~1 a la vía~1, de modo que la vía~1 contiene
        siempre la instrucción más joven.
  \item \textbf{Emitir} (\emph{issue}): pasar una instrucción de Decode a Execute. La
        \emph{Issue Unit} decide cuántas instrucciones del par se emiten; si solo emite el
        slot~0, el slot~1 se \emph{reagrupa}: pasa a ser el slot~0 del par siguiente
        (señal \sig{RegroupD}).
\end{itemize}

La Fig.~\ref{fig:base} muestra su datapath. La vía~0 (arriba) es el pipeline de cinco etapas
y la vía~1 (abajo) es su imagen especular. Sobre el eje central se ubican
los elementos que comparten ambas vías: la memoria de instrucciones, con dos puertos de
lectura; el registro IF/ID, que guarda el par completo; el register file, con cuatro
lecturas y dos escrituras; la Branch Unit y la memoria de datos, con un solo puerto. En la
parte inferior están la Issue Unit y la Hazard Unit. Las figuras del datapath usan cinco
colores: negro para los datos, azul para el control, naranja para la Hazard Unit y sus
señales, verde para la Issue Unit y morado para la Branch Unit. Un punto indica una unión
entre cables; dos cables que se cruzan sin punto no están conectados.

\begin{landscape}
\begin{figure}[p]
\centering
\includegraphics[width=\linewidth,height=0.88\textheight,keepaspectratio]{proc-dual}
\caption{Procesador dual-issue de partida: dos vías, Issue Unit con reagrupamiento, Branch
Unit compartida, stalls por \ins{lw} y flushes por salto en cualquier slot}
\label{fig:base}
\end{figure}
\end{landscape}

Los saltos (\ins{beq} y \ins{jal}) los resuelve un único recurso compartido por ambas vías,
la \emph{Branch Unit}, en la etapa Execute. La Branch Unit compara los operandos del
\ins{beq}, calcula la dirección de destino $\sig{PCTargetE} = \sig{PCE} + \sig{ImmExtE}$ y
activa \sig{PCSrcE} si el salto se toma. Mientras el salto avanza hasta Execute, Fetch sigue
buscando en secuencia; dicho de otro modo, el procesador predice que \emph{ningún} salto se
toma. Un \ins{jal} siempre se toma; un \ins{beq}, solo si sus operandos son iguales. Si el
salto se toma, la Hazard Unit descarta las instrucciones del camino incorrecto:
\begin{itemize}
  \item el par que está en Decode, con \sigk{FlushE}{0} y \sigk{FlushE}{1};
  \item el par que está en Fetch, con \sig{FlushD};
  \item si el salto está en el slot~0, también su compañero del slot~1, que es más joven,
        con \sigk{FlushM}{1}.
\end{itemize}

\paragraph{Diagramas de ejecución.} La ejecución se representa con diagramas como el de la
Fig.~\ref{fig:pd-nopred}. Cada fila corresponde a una instrucción y cada columna a un ciclo.
Cada celda indica el componente que la instrucción usa en ese ciclo: memoria de
instrucciones (IM), register file (RF), ALU, memoria de datos (DM) y, de nuevo, register file
en la escritura. Las celdas grises tachadas son instrucciones que se descartan; las flechas
azules representan forwarding, y la columna \textsf{v0}/\textsf{v1} indica la vía en la que
se ejecuta cada instrucción.

\subsection{El costo de un salto tomado}

La Fig.~\ref{fig:pd-nopred} muestra el bucle de la Tabla~\ref{tab:loop}, de tres pares
por iteración, con el \ins{beq} en el slot~0. En el ciclo~5, el \ins{beq} se resuelve como
tomado: se anulan su compañero más joven (\sigk{FlushM}{1}), el par que está en Decode y
el par que está en Fetch. La iteración siguiente empieza a buscarse recién en el ciclo~6.
Cada iteración ocupa cinco ciclos, aunque solo tres de ellos emiten instrucciones útiles.

\begin{table}[htbp]
\centering\small
\caption{Bucle de ejemplo: \ins{x1} cuenta de 1 a 10 y el \ins{beq} vuelve a \ins{loop}
mientras $\ins{x1} < 10$, así que el bucle da diez iteraciones}
\label{tab:loop}
\begin{tabular}{@{}rll@{}}
\toprule
Dirección & Instrucción & \\
\midrule
0--12 & \multicolumn{2}{l}{inicialización: $\ins{x1} = 0$, $\ins{x2} = 10$, $\ins{x3} = 1$, $\ins{x10} = 0$} \\
16 & \ins{addi x1, x1, 1} & \ins{loop}; par (16, 20) \\
20 & \ins{addi x10, x10, 3} & \\
24 & \ins{slt  x4, x1, x2} & par (24, 28): $\ins{x4} = (\ins{x1} < \ins{x2})$ \\
28 & \ins{addi x0, x0, 0} & \\
32 & \ins{beq  x4, x3, loop} & par (32, 36): salto en el slot 0 \\
36 & \ins{addi x5, x0, 7} & se ejecuta al salir del bucle \\
\bottomrule
\end{tabular}
\end{table}

\begin{figure}[htbp]
\centering
\includegraphics[width=0.92\linewidth]{pd-nopred}
\caption{Dual-issue sin predicción: el \ins{beq} tomado descarta dos pares y su compañero
del slot~1. La iteración dura cinco ciclos}
\label{fig:pd-nopred}
\end{figure}

La pérdida es fácil de cuantificar: si un programa ejecuta $N_T$ saltos tomados, el
procesador pierde $2N_T$ ciclos (los \ins{jal} cuentan como saltos tomados). Estos ciclos
se pierden porque Fetch no sabe a dónde ir
hasta que el salto llega a Execute. Para evitarlo, Fetch tendría que decidir de dónde buscar
las instrucciones siguientes \emph{antes} de conocer el resultado del salto, es decir,
tendría que \emph{predecirlo}. Esta técnica es la predicción de saltos, que se presenta
en la sección siguiente.

% =====================================================================
\section{Visión general}
\label{sec:overview}

Un \emph{predictor de saltos} adivina, en Fetch, si un salto se va a tomar y, si se toma, a
qué dirección va. Si acierta, el pipeline sigue lleno y el salto no cuesta nada. Si se
equivoca, el error se detecta en Execute, como ahora, y se pagan los mismos dos ciclos que
hoy cuesta cada salto tomado. Visto así, nuestro procesador ya tiene un predictor: el más
simple de todos, que siempre dice ``no tomado''.

Ese predictor es \emph{estático}: su respuesta no depende de lo que haga el programa. Un
predictor \emph{dinámico}, en cambio, recuerda lo que hizo cada salto las veces anteriores
y predice a partir de esa historia. En~\cite{harris}, el predictor dinámico es una máquina
de estados de 2 bits por salto, que repasamos en la sección~\ref{sec:counter}. Esa máquina
decide si el salto se toma, pero por sí sola no alcanza para predecir en Fetch. La
sección~\ref{sec:btb} muestra qué le falta y cómo se resuelve: guardando los contadores,
junto con el destino de cada salto, en una tabla que Fetch consulta con el PC. Esa tabla es
el \emph{Branch Target Buffer} (BTB).

\paragraph{Qué hace falta.} El predictor necesita cuatro piezas de hardware, que
desarrollamos en las secciones siguientes:
\begin{enumerate}
  \item un \textbf{contador de 2 bits} por salto, que decide la dirección
        (sección~\ref{sec:counter});
  \item la \textbf{tabla BTB}, que guarda el contador y el destino de cada salto y que
        Fetch lee con el PC, una vez por cada instrucción del par (sección~\ref{sec:btb});
  \item una etapa \textbf{Fetch} que use esa predicción para elegir el siguiente PC y que
        anule las instrucciones que quedan fuera del camino predicho
        (sección~\ref{sec:fetch});
  \item una \textbf{Branch Unit} que, en Execute, compare la predicción con el resultado
        real, corrija el PC si hubo error y actualice la tabla
        (secciones~\ref{sec:carry} y~\ref{sec:verify}).
\end{enumerate}
El BTB queda dentro de la etapa Fetch, y el resto de los cambios afecta a la Branch Unit, a
los registros de pipeline y a dos señales de la Hazard Unit. La Issue Unit, el forwarding,
el register file y la memoria de datos no cambian. En las figuras, los elementos agregados
para la predicción aparecen en café; el resto conserva los colores de la
Fig.~\ref{fig:base}.

% =====================================================================
\section{El contador de 2 bits}
\label{sec:counter}

El predictor dinámico más simple usa 1 bit por salto: recuerda si el salto se tomó la
última vez y predice que hará lo mismo. Consideremos el \ins{beq} que cierra un bucle de
cuatro iteraciones. Se toma tres veces, para volver al inicio, y no se toma la cuarta, cuando
el bucle termina. Si el bucle se ejecuta dos veces, el salto produce la secuencia
\textsf{T T T N T T T N} (T = tomado, N = no tomado), y el predictor de 1 bit se equivoca en
la cuarta ejecución, porque el bit dice ``tomado'', y también en la quinta: al salir del bucle
el bit quedó en ``no tomado'' y la primera iteración de la ejecución siguiente vuelve a
predecirse mal. En régimen, son \emph{dos} errores por cada ejecución del bucle
(Tabla~\ref{tab:1v2}).

El contador de 2 bits resuelve este problema con cuatro estados
(Fig.~\ref{fig:fsm}). Los dos estados de la izquierda (\texttt{11}, \texttt{10})
predicen tomado y los dos de la derecha (\texttt{01}, \texttt{00}) predicen no tomado; la
predicción es simplemente el bit más significativo, $\sig{Ctr}[1]$. Cada salto tomado
mueve el contador hacia \texttt{11} y cada salto no tomado lo mueve hacia \texttt{00}, y en
los extremos el contador satura. Un solo resultado contrario, como la salida de un bucle,
lleva el contador de \texttt{11} a \texttt{10}, que todavía predice tomado. Por eso el
contador de 2 bits se equivoca solo en la última iteración de cada bucle: en la
Tabla~\ref{tab:1v2}, a partir de la segunda ejecución del bucle, comete un error en vez de
dos.

\begin{table}[htbp]
\centering\small
\caption{Predicciones del \ins{beq} que cierra un bucle de cuatro iteraciones ejecutado dos
veces. Los errores están marcados con $\times$. La primera vez, ningún predictor conoce el
salto y ambos dicen ``no tomado''}
\label{tab:1v2}
\begin{tabular}{@{}lcccccccc@{}}
\toprule
Ejecución del \ins{beq} & 1 & 2 & 3 & 4 & 5 & 6 & 7 & 8 \\
Resultado real          & T & T & T & N & T & T & T & N \\
\midrule
1 bit: predicción       & N $\times$ & T & T & T $\times$ & N $\times$ & T & T & T $\times$ \\
2 bits: contador antes  & --- & \texttt{10} & \texttt{11} & \texttt{11} & \texttt{10} & \texttt{11} & \texttt{11} & \texttt{11} \\
2 bits: predicción      & N $\times$ & T & T & T $\times$ & T & T & T & T $\times$ \\
\bottomrule
\end{tabular}
\end{table}

\begin{figure}[htbp]
\centering
\includegraphics[width=\linewidth]{fsm-2bit}
\caption{Contador de 2 bits (adaptado de~\cite{harris}). En café, el valor inicial del contador}
\label{fig:fsm}
\end{figure}

Falta decidir con qué valor empieza el contador. Mientras un salto no se ha tomado nunca,
el predictor no sabe nada de él y predice ``no tomado''. La primera vez que un \ins{beq} se
toma, su contador empieza en \texttt{10} (\emph{weakly taken}): el salto se acaba de tomar,
pero un solo resultado contrario basta para cambiar la predicción. Un \ins{jal} se toma
siempre, así que su contador empieza directamente en \texttt{11}.

% =====================================================================
\section{El Branch Target Buffer}
\label{sec:btb}

\subsection{Por qué hace falta una tabla}

Hasta aquí, el predictor es una máquina de estados con un registro de 2 bits. Para usarla en
Fetch hay que responder tres preguntas, y en Fetch no se conoce ninguna de las respuestas:
\begin{itemize}
  \item \textbf{¿La instrucción es un salto?} En Fetch, la instrucción acaba de salir de la
        memoria de instrucciones y todavía no se ha decodificado. Eso ocurre recién en
        Decode, un ciclo después.
  \item \textbf{¿Qué contador le corresponde?} Un programa tiene muchos saltos, y cada uno
        necesita su propio contador: el \ins{beq} de un bucle no se comporta igual que el
        de otro. Lo único que Fetch sabe de la instrucción es su dirección, \sig{PCF}, así
        que el contador tiene que encontrarse a partir de ella.
  \item \textbf{Si se toma, ¿a dónde va?} El destino, $\sig{PC} + \sig{imm}$, lo calcula
        la Branch Unit en Execute. Si Fetch predice ``tomado'', necesita ese destino en
        ese mismo ciclo para buscar la instrucción siguiente.
\end{itemize}
Las tres se resuelven con una misma memoria pequeña, que Fetch lee con el PC. Cada entrada
de la memoria recuerda un salto que ya se ejecutó: su contador de 2 bits y la dirección a
la que fue la última vez que se tomó. Si el PC buscado tiene una entrada, la instrucción es
un salto conocido, el contador dice si predecir tomado y la entrada entrega el destino. Si no
tiene, se predice ``no tomado'', igual que ahora. Esta memoria es el \emph{Branch Target
Buffer}. En~\cite{harris} aparece en una sola frase, como una tabla de los últimos saltos
ejecutados que guarda su destino y si se tomaron; en esta sección la construimos.

Observe que la tabla no guarda instrucciones, solo direcciones. Cuando predice ``tomado'',
el siguiente PC pasa a ser el destino guardado y, en el ciclo siguiente, Fetch busca en la
memoria de instrucciones a partir de esa dirección, como con cualquier otro PC.

\subsection{Elegir la entrada: el índice}

Nuestra tabla tiene 16 entradas, numeradas de 0 a~15. Para que la lectura sea rápida, no se
recorren las 16: cada dirección tiene \emph{un único} lugar posible en la tabla, y ese
lugar se obtiene directamente de los bits del PC. Esta organización se llama de
\emph{correspondencia directa}.

Para numerar 16 entradas hacen falta 4~bits. Los dos bits más bajos del PC no sirven,
porque las instrucciones ocupan 4~bytes y sus direcciones son siempre múltiplos de~4:
$\sig{PC}[1{:}0] = \texttt{00}$. Se usan los cuatro bits siguientes, $\sig{PC}[5{:}2]$, que
llamamos el \emph{índice}. La Tabla~\ref{tab:idx} lo calcula para las instrucciones del
bucle de la Tabla~\ref{tab:loop}. Por ejemplo, $32 = \texttt{10\,0000}$ en binario; sus
bits 5 a~2 son \texttt{1000}, es decir, la entrada~8. Como el índice son los bits que siguen
a los dos ceros, instrucciones consecutivas caen en entradas consecutivas.

\begin{table}[htbp]
\centering\small
\caption{Índice de las instrucciones del bucle de la Tabla~\ref{tab:loop} (6 bits bajos del
PC)}
\label{tab:idx}
\begin{tabular}{@{}rlccc@{}}
\toprule
PC & Instrucción & $\sig{PC}[5{:}0]$ & $\sig{PC}[5{:}2]$ & Entrada \\
\midrule
16 & \ins{addi x1, x1, 1}   & \texttt{01\,00\,00} & \texttt{0100} & 4 \\
20 & \ins{addi x10, x10, 3} & \texttt{01\,01\,00} & \texttt{0101} & 5 \\
24 & \ins{slt  x4, x1, x2}  & \texttt{01\,10\,00} & \texttt{0110} & 6 \\
28 & \ins{addi x0, x0, 0}   & \texttt{01\,11\,00} & \texttt{0111} & 7 \\
32 & \ins{beq  x4, x3, loop}& \texttt{10\,00\,00} & \texttt{1000} & 8 \\
36 & \ins{addi x5, x0, 7}   & \texttt{10\,01\,00} & \texttt{1001} & 9 \\
\bottomrule
\end{tabular}
\end{table}

Solo el \ins{beq} llega a escribirse en la tabla, en la entrada~8. Las demás instrucciones
también leen su entrada en Fetch, pero la encuentran vacía y se predicen ``no tomado''.

\subsection{Reconocer la instrucción: Tag y bit de validez}

Con 16 entradas y muchas más direcciones, varias direcciones comparten el mismo índice. La
dirección~96, por ejemplo, es $\texttt{1\,1000\,00}$: sus bits 5 a~2 también son
\texttt{1000}. Si el \ins{beq} de la dirección~32 está guardado en la entrada~8 y Fetch busca
la dirección~96, el índice por sí solo no permite distinguirlas, y la instrucción de la
dirección~96 recibiría la predicción del \ins{beq}.

Para distinguirlas, cada entrada guarda también el resto del PC, los 26 bits
$\sig{PC}[31{:}6]$, en un campo llamado \emph{Tag}. Para la dirección~32 el Tag vale~0, y
para la~96 vale~1. Al leer la tabla, se compara el Tag guardado con los bits
$\sig{PC}[31{:}6]$ de la dirección buscada. Como el Tag y el índice juntos forman el PC
completo, si coinciden, la entrada corresponde \emph{exactamente} a esa instrucción.

Falta un último detalle. Después del reset, todas las entradas tienen el Tag en~0, y la
dirección~32 coincidiría con la entrada~8 aunque nunca se haya escrito nada. Por eso cada
entrada tiene un bit de validez, V, que el reset pone en~0 y que se activa cuando se escribe
la entrada.

Así, cada entrada tiene cuatro campos: V (1~bit), Tag (26~bits), el contador Ctr (2~bits) y
el destino Target (32~bits). Leer la tabla con una dirección tiene dos resultados posibles:
\begin{itemize}
  \item \textbf{acierto} (\emph{hit}): la entrada es válida y el Tag coincide. La
        instrucción es un salto conocido; se predice tomado si el contador está en
        \texttt{11} o \texttt{10}, es decir, si $\sig{Ctr}[1] = 1$, y el destino es Target;
  \item \textbf{fallo} (\emph{miss}): en cualquier otro caso. Se predice no tomado.
\end{itemize}
En una sola ecuación,
\[
  \sig{predtaken} = \sig{V}\cdot(\sig{Tag} = \sig{PC}[31{:}6])\cdot\sig{Ctr}[1].
\]
El BTB no necesita saber si la instrucción es un salto. Como solo se escriben saltos (lo
vemos a continuación), una dirección con entrada válida es siempre la de un salto. Eso
incluye al \ins{jal}: la primera vez falla en la tabla y paga los dos ciclos de siempre,
pero deja su entrada con el contador en \texttt{11}, y como siempre se toma, las veces
siguientes se predice bien.

\subsection{Escribir la tabla}

La tabla se escribe en un solo lugar: la Branch Unit, en Execute, cada vez que hay un salto
en esa etapa, porque recién ahí se sabe si se tomó y cuál es su destino. Con el índice del
PC del salto lee la entrada que le corresponde, comprueba si hay hit, y en el flanco de
subida del reloj escribe el resultado:
\begin{itemize}
  \item \textbf{salto tomado}: se escriben V~=~1, el Tag y el destino calculado. El contador
        aumenta en~1 si ya había hit; si la entrada es nueva, empieza en \texttt{10}
        (\texttt{11} para un \ins{jal}), como se vio en la sección~\ref{sec:counter};
  \item \textbf{salto no tomado con hit}: el contador disminuye en~1;
  \item \textbf{salto no tomado sin hit}: no se escribe nada. Ese salto se predijo ``no
        tomado'' y acertó, así que no hay nada que recordar.
\end{itemize}
Si la entrada estaba ocupada por otro salto con el mismo índice, el salto tomado la
reemplaza. Este conflicto nunca produce un resultado incorrecto: a lo sumo, el salto que se
perdió se predice ``no tomado'' la próxima vez, y si eso es un error, la verificación de la
sección~\ref{sec:verify} lo corrige.

La Tabla~\ref{tab:trace} sigue la entrada~8 durante las diez iteraciones del bucle de la
Tabla~\ref{tab:loop}. En la primera iteración la entrada está vacía: hay fallo, se predice
``no tomado'' y el \ins{beq} paga los dos ciclos de siempre. Al resolverse, se crea la
entrada. Desde la segunda iteración hay hit y la predicción acierta, hasta la última, en la
que el \ins{beq} no se toma y el contador baja a~\texttt{10}. Si el bucle se vuelve a
ejecutar, la entrada sigue ahí y su primer salto se predice bien: el único error es el de la
salida.

\begin{table}[htbp]
\centering\small
\caption{Entrada~8 del BTB durante el bucle de la Tabla~\ref{tab:loop}. Los errores de
predicción están marcados con $\times$}
\label{tab:trace}
\begin{tabular}{@{}cccc c>{\raggedright\arraybackslash}p{3.9cm}@{}}
\toprule
Iteración & Entrada~8 antes & Lectura & Predicción & Resultado & Entrada~8 después \\
\midrule
1        & V = 0              & fallo & no tomado $\times$ & tomado    & V = 1, Tag = 0, Ctr = \texttt{10}, Target = 16 \\
2        & Ctr = \texttt{10}  & hit   & tomado             & tomado    & Ctr = \texttt{11} \\
3 a 9    & Ctr = \texttt{11}  & hit   & tomado             & tomado    & Ctr = \texttt{11} (satura) \\
10       & Ctr = \texttt{11}  & hit   & tomado $\times$    & no tomado & Ctr = \texttt{10} \\
\bottomrule
\end{tabular}
\end{table}

Es importante que la escritura parta del contador que hay \emph{en ese momento} en la tabla,
y no del valor que se leyó en Fetch cuando se predijo el salto. La predicción y la escritura
están separadas por dos ciclos, y en un bucle corto la instancia siguiente del mismo salto se
busca antes de que la anterior se escriba. Supongamos un bucle de dos pares: el \ins{beq} se
busca en el ciclo~$t$, llega a Execute en el ciclo~$t+2$ y, en ese mismo ciclo, Fetch ya
vuelve a buscarlo para la iteración siguiente. Si cada escritura partiera del valor leído en
Fetch, las dos instancias habrían leído el mismo contador y escribirían el mismo resultado:
se perdería una cuenta. Leyendo la tabla en Execute, cada escritura parte del valor que dejó
la anterior.

\subsection{Dos puertos de lectura}

El procesador busca dos instrucciones por ciclo, y cualquiera de ellas puede ser un salto,
así que el BTB se lee dos veces por ciclo: tiene dos puertos de lectura, del mismo modo que
la memoria de instrucciones. El puerto~0 recibe \sig{PCF}, la dirección del slot~0, y
entrega \sigk{PredF}{0} (se predice tomado) y \sigk{TargetF}{0} (destino predicho). El
puerto~1 recibe \sig{PCPlus4F}, la dirección del slot~1, y entrega \sigk{PredF}{1} y
\sigk{TargetF}{1}. Como las dos direcciones son consecutivas, los dos puertos leen
entradas vecinas: el par (32,~36) lee las entradas~8 y~9. La lectura es combinacional,
igual que la de la memoria de instrucciones, de modo que la predicción está lista en el
mismo ciclo en que se busca el par.

La Fig.~\ref{fig:btb} reúne todo lo anterior. En el centro está la tabla, con el contenido
de la entrada~8 durante el bucle. Arriba está el puerto~0: el PC se divide en Tag, índice y
los dos ceros; el índice elige una fila (el multiplexor 16:1), un comparador verifica el
Tag, y una compuerta AND combina V, el resultado del comparador y $\sig{Ctr}[1]$ para dar
\sigk{PredF}{0}. Abajo está el puerto~1, idéntico. A la derecha, el bloque de escritura
recibe de la Branch Unit las señales del salto que está en Execute; los nombres del módulo
(\sig{pc\_0}, \sig{upc}, \dots) aparecen junto a la señal del datapath a la que se conectan:
\begin{itemize}
  \item \sig{update} $\leftarrow$ $\sig{CtrlE} = \sig{BranchE} + \sig{JumpE}$: en Execute hay
        un \ins{beq} o un \ins{jal};
  \item \sig{upc} $\leftarrow$ \sig{PCE}, la dirección del salto, que da el índice y el Tag;
  \item \sig{utaken} $\leftarrow$ \sig{PCSrcE}, si el salto se tomó;
  \item \sig{ujump} $\leftarrow$ \sig{JumpE}, si es un \ins{jal};
  \item \sig{utarget} $\leftarrow$ \sig{PCTargetE}, el destino calculado.
\end{itemize}

\begin{figure}[htbp]
\centering
\includegraphics[width=\linewidth]{btb-inside}
\caption{Estructura interna del BTB: dos puertos de lectura (uno por slot) y el bloque de
escritura. El contenido de ejemplo es la entrada~8 durante el bucle de la
Tabla~\ref{tab:loop}}
\label{fig:btb}
\end{figure}

\begin{hdl}{Lectura y actualización del BTB}
// prediction (Fetch): one read port per slot, combinational
wire [3:0] idx_0 = pc_0[5:2];
wire       hit_0 = valid[idx_0] & (tag[idx_0] == pc_0[31:6]);
assign predtaken_0 = hit_0 & ctr[idx_0][1];
assign target_0    = target[idx_0];
// (slot 1: the same with pc_1)

// update (Execute): reads the current entry, writes it at the clock edge
wire [3:0] uidx = upc[5:2];
wire       uhit = valid[uidx] & (tag[uidx] == upc[31:6]);
always @(posedge clk or posedge reset)
  if (reset) /* every valid <= 0 */ ;
  else if (update & utaken) begin            // taken: write the entry
    valid[uidx]  <= 1'b1;
    tag[uidx]    <= upc[31:6];
    target[uidx] <= utarget;
    ctr[uidx]    <= ujump ? 2'b11 : uhit ? inc : 2'b10;
  end else if (update & uhit)                // not taken: count down
    ctr[uidx]    <= dec;                     // inc / dec saturate at 11 / 00
\end{hdl}

\subsection{Resumen}

La Tabla~\ref{tab:params} reúne las decisiones de diseño del BTB que se explicaron en esta
sección y en la anterior.

\begin{table}[htbp]
\centering\small
\caption{Parámetros del predictor}
\label{tab:params}
\begin{tabular}{@{}l>{\raggedright\arraybackslash}p{5.0cm}>{\raggedright\arraybackslash}p{5.6cm}@{}}
\toprule
Parámetro & Valor & Motivo \\
\midrule
Entradas & 16, correspondencia directa & los saltos de un bucle pequeño caben sin compartir entrada \\
Contenido de una entrada & V, Tag (26 b), Ctr (2 b), Target (32 b) & 61 bits; 976 bits en total \\
Puertos & 2 lecturas, 1 escritura & una predicción por slot; un salto por ciclo \\
Fallo en la tabla & se predice no tomado & igual que el dual-issue de partida \\
Entrada nueva & solo si el salto se toma & un salto nunca tomado no necesita entrada \\
Verificación & en Execute, en la Branch Unit & el salto se resuelve donde se resolvía (sección~\ref{sec:verify}) \\
Penalización por error & 2 ciclos & la misma que un salto tomado sin predicción (sección~\ref{sec:verify}) \\
\bottomrule
\end{tabular}
\end{table}

% =====================================================================
\section{Fetch con predicción}
\label{sec:fetch}

\subsection{El siguiente PC}

En el dual-issue de partida, el siguiente PC era $\sig{PC}+8$ (se emitió el par),
$\sig{PC}+4$ (reagrupamiento) o \sig{PCTargetE} (salto tomado). Con predicción aparecen dos
fuentes nuevas: los destinos predichos \sigk{TargetF}{0} y \sigk{TargetF}{1}, que entrega
el BTB. Además, \sig{PCTargetE} se reemplaza por \sig{PCCorrectE}. Esta señal, junto con
\sig{MispredictE}, la genera la Branch Unit (sección~\ref{sec:verify}): \sig{MispredictE}
indica que la predicción del salto que está en Execute fue incorrecta, y \sig{PCCorrectE}
es la dirección por la que Fetch debe continuar. Los cinco candidatos se eligen en este
orden de prioridad:
\[
  \sig{PCNextF} =
  \begin{cases}
    \sig{PCCorrectE} & \text{si } \sig{MispredictE} \\
    \sigk{TargetF}{0} & \text{si } \sigk{PredF}{0} \\
    \sig{PCPlus4F}   & \text{si } \sig{RegroupD} \\
    \sigk{TargetF}{1} & \text{si } \sigk{PredF}{1} \\
    \sig{PCPlus8F}   & \text{en otro caso}
  \end{cases}
\]
y se implementan con una cadena de cuatro multiplexores 2:1 (Fig.~\ref{fig:zoom-fetch}). El
orden no es arbitrario:
\begin{itemize}
  \item \textbf{\sig{MispredictE} primero.} Si la Branch Unit detecta un error, todo lo que
        hay en Fetch y Decode pertenece al camino incorrecto, incluidas sus predicciones.
  \item \textbf{\sigk{PredF}{0} antes que \sig{RegroupD}.} Recordemos que en un
        reagrupamiento el nuevo par se forma con el slot~1 que espera en Decode y con
        \sigk{InstrF}{0}, y que el PC avanza solo en~4, hasta la instrucción que sigue a
        \sigk{InstrF}{0}. Pero si \sigk{InstrF}{0} es un salto predicho tomado, la
        instrucción que la sigue en el camino predicho es su destino, no
        $\sig{PCF}+4$.
  \item \textbf{\sig{RegroupD} antes que \sigk{PredF}{1}.} Durante un reagrupamiento,
        \sigk{InstrF}{1} no entra a IF/ID: se descarta y se vuelve a buscar en el ciclo
        siguiente. Su predicción, por lo tanto, todavía no corresponde.
\end{itemize}

\subsection{Anulación del slot 1}

Si el slot~0 del par buscado es un salto predicho tomado, el slot~1 es la instrucción que
lo sigue en memoria, que queda \emph{fuera} del camino predicho. No puede ejecutarse, así
que se anula antes de entrar a IF/ID:
\[
  \sigk{AnnulF}{1} = \sigk{PredF}{0}\cdot\overline{\sig{RegroupD}}.
\]
Con \sigk{AnnulF}{1} activa, dos multiplexores ponen en~0 la instrucción del slot~1 y su
predicción. Una instrucción en~0 actúa como burbuja: no escribe registros ni memoria y no
es un salto. Durante un reagrupamiento no se anula nada, porque el slot~1 nuevo es
\sigk{InstrF}{0} y \sigk{InstrF}{1} ya se descarta.

Esta anulación tiene un costo: cuando el salto ocupa el slot~0, el par sale de Fetch con una
sola instrucción útil. Si el salto ocupa el slot~1, en cambio, el par completo es útil,
porque el slot~0 es más antiguo que el salto.

\subsection{Reagrupamiento de la predicción}

Si la instrucción del slot~1 queda en espera y se reagrupa, debe llevarse su predicción. Por
eso se agregan dos multiplexores de reagrupamiento más, con la misma lógica que los de la
instrucción y su PC:
\begin{itemize}
  \item \textbf{nuevo slot~0} $\leftarrow$ \{\sigk{PredD}{1}, \sigk{TargetD}{1}\},
        realimentado desde la salida de IF/ID;
  \item \textbf{nuevo slot~1} $\leftarrow$ \{\sigk{PredF}{0}, \sigk{TargetF}{0}\}, la
        predicción de \sigk{InstrF}{0}.
\end{itemize}

El reagrupamiento sigue funcionando sin otros cambios. En el dual-issue de partida, mientras
un par está en Decode, Fetch se encuentra en $\sig{PCF} = \sigk{PCD}{1} + 4$, de modo que
\sigk{InstrF}{0} es la instrucción que sigue al slot~1, y por eso puede completar el par
reagrupado. Con predicción, \sig{PCF} apunta a la instrucción que sigue
al slot~1 \emph{en el camino predicho}: si el slot~1 es un salto predicho tomado, Fetch ya
está buscando su destino.

Por ejemplo, supongamos que en Decode está el par formado por \ins{slt x4, x1, x2} en el
slot~0 y \ins{beq x4, x3, loop} en el slot~1, y que el \ins{beq} se predijo tomado. El \ins{beq} necesita el \ins{x4} que calcula
el \ins{slt} en el mismo ciclo, así que la Issue Unit emite solo el \ins{slt} y el
\ins{beq} se reagrupa. Mientras tanto, como el \ins{beq} se predijo tomado, Fetch ya está
buscando la primera instrucción de \ins{loop}. El nuevo par queda formado por el \ins{beq}
y esa instrucción, que es exactamente la que lo sigue en el camino predicho. El
reagrupamiento no necesita saber que hubo una predicción: solo toma lo que Fetch le
entrega.

\begin{figure}[htbp]
\centering
\includegraphics[height=13.0cm]{proc-btb-zoom-fetch}
\caption{Fetch con predicción: BTB, cadena de multiplexores del siguiente PC,
reagrupamiento de \{Pred, Target\} y anulación del slot~1 con \sigk{AnnulF}{1}}
\label{fig:zoom-fetch}
\end{figure}

\begin{hdl}{Siguiente PC, anulación y reagrupamiento de la predicción}
// priority: misprediction, slot 0 predicted taken, regroup, slot 1 predicted taken
assign PCNextF = MispredictE ? PCCorrectE :
                 PredF_0     ? TargetF_0  :
                 RegroupD    ? PCPlus4F   :
                 PredF_1     ? TargetF_1  :
                               PCPlus8F;

// slot 0 predicted taken: the fetched slot 1 is not on the predicted path
wire AnnulF_1 = PredF_0 & ~RegroupD;

wire [31:0] InstrN_1 = RegroupD ? InstrF_0 : AnnulF_1 ? 32'b0 : InstrF_1;
wire [32:0] PredN_0  = RegroupD ? {PredD_1, TargetD_1} : {PredF_0, TargetF_0};
wire [32:0] PredN_1  = RegroupD ? {PredF_0, TargetF_0} : AnnulF_1 ? 33'b0 : {PredF_1, TargetF_1};
\end{hdl}

% =====================================================================
\section{Transporte de la predicción}
\label{sec:carry}

La predicción se hace en Fetch, pero se verifica en Execute, dos ciclos después. Hasta
entonces viaja con la instrucción, en los mismos registros de pipeline: 33 bits por
instrucción (\sig{Pred} y \sig{Target}). La Tabla~\ref{tab:regs} muestra los registros que
crecen. EX/MEM y MEM/WB no cambian, porque la predicción ya no se usa después de Execute.

\begin{table}[htbp]
\centering\small
\caption{Registros de pipeline que crecen ($k = 0, 1$)}
\label{tab:regs}
\begin{tabular}{@{}lll@{}}
\toprule
Registro & Campos nuevos & Bits agregados \\
\midrule
IF/ID (uno para el par) & \sigk{PredD}{0}, \sigk{TargetD}{0}, \sigk{PredD}{1}, \sigk{TargetD}{1} & 66 \\
ID/EX$_k$ & \sigk{PredE}{k}, \sigk{TargetE}{k} & 33 \\
\bottomrule
\end{tabular}
\end{table}

Como viajan en los mismos registros, la predicción obedece a los mismos stalls y flushes
que la instrucción: si la instrucción se detiene, su predicción también; si se descarta,
su predicción desaparece con ella. Por eso basta con que la Branch Unit mire la
predicción de la vía que contiene el salto.

% =====================================================================
\section{Verificación en la Branch Unit}
\label{sec:verify}

\subsection{Error de predicción}

La Branch Unit sigue resolviendo el salto como en el dual-issue de partida: selecciona los operandos
de la vía que contiene el salto ($\sig{BrLaneE} = \sigk{BranchE}{1} + \sigk{JumpE}{1}$),
compara $\sig{SrcA} = \sig{SrcB}$ y calcula $\sig{PCTargetE} = \sig{PCE} + \sig{ImmExtE}$
(Fig.~\ref{fig:bu}). Ahora, además, toma con los mismos multiplexores la predicción de esa
vía, \sig{PredTakenE} y \sig{PredTargetE}, y la compara con el resultado. La predicción es
incorrecta si se equivocó en la dirección o, cuando el salto se toma, en el destino:
\[
  \sig{MispredictE} = \sig{CtrlE}\cdot\big[(\sig{PCSrcE} \neq \sig{PredTakenE}) +
  \sig{PCSrcE}\cdot(\sig{PCTargetE} \neq \sig{PredTargetE})\big],
\]
donde $\sig{CtrlE} = \sig{BranchE} + \sig{JumpE}$ indica que hay un salto en Execute. Con
el Tag completo, un hit implica el mismo PC y, por lo tanto, el mismo destino, así que la
comparación del destino nunca falla en este diseño. La conservamos porque hace que la
verificación sea correcta por sí misma, aunque la tabla se simplificara (con un Tag más
corto, por ejemplo).

La Tabla~\ref{tab:cases} enumera los casos. Lo importante es que \sig{MispredictE}
reemplaza a \sig{PCSrcE} como señal de ``el camino es incorrecto''. Un salto tomado que se
predijo como tomado no genera ningún flush, y un salto predicho como tomado que no se toma
sí lo genera.

\begin{table}[htbp]
\centering\small
\caption{Casos de la verificación en Execute}
\label{tab:cases}
\begin{tabular}{@{}llll@{}}
\toprule
Predicción & Resultado & \sig{MispredictE} & Siguiente PC \\
\midrule
no tomado & no tomado & 0 & sigue la secuencia \\
no tomado & tomado    & 1 & $\sig{PCCorrectE} = \sig{PCTargetE}$ \\
tomado    & tomado    & 0 & Fetch ya está en el destino \\
tomado    & no tomado & 1 & $\sig{PCCorrectE} = \sig{PCPlus4E}$ \\
\bottomrule
\end{tabular}
\end{table}

\subsection{Corrección del PC}

Cuando la predicción falla, Fetch debe volver al camino correcto: al destino si el salto se
tomó, o a la instrucción siguiente al salto si no se tomó:
\[
  \sig{PCCorrectE} = \sig{PCSrcE}\;?\;\sig{PCTargetE} : \sig{PCPlus4E}.
\]
Por eso la Branch Unit recibe ahora también el $\sig{PC}+4$ de cada vía. En el dual-issue
de partida este caso no existía: un salto no tomado nunca desviaba a Fetch. Observe que,
cuando \sigk{AnnulF}{1} anuló en Fetch el slot~1 que seguía al salto, la dirección de esa
instrucción es justamente \sig{PCPlus4E}, de modo que la instrucción anulada se vuelve a
buscar.

\subsection{Anulación de la otra vía}

En el dual-issue de partida, $\sigk{TakenE}{0} = \sig{PCSrcE}\cdot\overline{\sig{BrLaneE}}$ indicaba un salto tomado en la vía~0 y anulaba
a su compañero más joven de la vía~1. Ahora la condición es un error de predicción en la
vía~0:
\[
  \sig{FlushLane1E} = \sig{MispredictE}\cdot\overline{\sig{BrLaneE}}.
\]
En ambos sentidos del error, la vía~1 contiene una instrucción del camino incorrecto:
\begin{itemize}
  \item si el salto se predijo no tomado y se tomó, la vía~1 contiene la instrucción
        siguiente en memoria;
  \item si se predijo tomado y no se tomó, la vía~1 contiene la primera instrucción del
        destino predicho. Esto ocurre cuando el salto llegó al slot~0 por un
        reagrupamiento. Si, en cambio, el salto ya estaba en el slot~0 cuando se buscó,
        \sigk{AnnulF}{1} dejó la vía~1 vacía y anularla de nuevo no tiene efecto.
\end{itemize}

\subsection{Actualización de la tabla}

Por último, la Branch Unit entrega al BTB las señales del puerto de actualización
(sección~\ref{sec:btb}): \sig{CtrlE}, el PC del salto (\sig{PCE}), \sig{PCSrcE},
\sig{JumpE} y \sig{PCTargetE}. Todas salen de los mismos multiplexores de selección de
vía, así que la actualización corresponde siempre al único salto que hay en Execute.

\begin{figure}[htbp]
\centering
\includegraphics[width=\linewidth]{bu-inside}
\caption{Branch Unit con verificación: selección de la vía, resolución, detección del error
de predicción, \sig{PCCorrectE}, \sig{FlushLane1E} y señales de actualización del BTB}
\label{fig:bu}
\end{figure}

\begin{hdl}{Verificación de la predicción}
wire BrLaneE = BranchE_1 | JumpE_1;          // lane that owns the branch unit
// ... lane multiplexers: PCE, PCPlus4E, ImmExtE, SrcAE, SrcBE, PredTakenE, PredTargetE
assign CtrlE  = BranchE | JumpE;
assign EqualE = (SrcAE == SrcBE);
adder pcaddbranch(PCE, ImmExtE, PCTargetE);
assign PCSrcE = (BranchE & EqualE) | JumpE;

// wrong direction, or taken to a different target
assign MispredictE = CtrlE & ((PCSrcE != PredTakenE) | (PCSrcE & (PCTargetE != PredTargetE)));
mux2 #(32) correctmux(PCPlus4E, PCTargetE, PCSrcE, PCCorrectE);
// mispredicted branch in lane 0: lane 1 is younger and on the wrong path
assign FlushLane1E = MispredictE & ~BrLaneE;
\end{hdl}

% =====================================================================
\section{Cambios en la Hazard Unit}
\label{sec:hazard}

La Hazard Unit no cambia de estructura: solo cambia la señal que indica que el camino es
incorrecto. Donde antes usaba \sig{PCSrcE} (salto tomado) ahora usa \sig{MispredictE}
(predicción incorrecta), y \sigk{TakenE}{0} se reemplaza por \sig{FlushLane1E}
(Tabla~\ref{tab:hu}). Las prioridades se mantienen: un error de predicción cancela el stall
por load-use y el reagrupamiento, porque el par que está en Decode se va a descartar.

\begin{table}[htbp]
\centering\small
\caption{Salidas de la Hazard Unit que cambian: dual-issue y dual-issue con predicción}
\label{tab:hu}
\renewcommand{\arraystretch}{1.15}
\begin{tabular}{@{}lll@{}}
\toprule
Salida & Dual-issue & Con predicción \\
\midrule
\sig{lwStallD} & $\sigk{LoadUseD}{0}\cdot\overline{\sig{PCSrcE}}$ & $\sigk{LoadUseD}{0}\cdot\overline{\sig{MispredictE}}$ \\
\sig{RegroupD} & $\sig{IssueOneD}\cdot\overline{\sig{lwStallD}}\cdot\overline{\sig{PCSrcE}}$ & $\sig{IssueOneD}\cdot\overline{\sig{lwStallD}}\cdot\overline{\sig{MispredictE}}$ \\
\sig{FlushD} & \sig{PCSrcE} & \sig{MispredictE} \\
\sigk{FlushE}{0} & $\sig{lwStallD} + \sig{PCSrcE}$ & $\sig{lwStallD} + \sig{MispredictE}$ \\
\sigk{FlushE}{1} & $\sig{lwStallD} + \sig{PCSrcE} + \sig{IssueOneD}$ & $\sig{lwStallD} + \sig{MispredictE} + \sig{IssueOneD}$ \\
\sigk{FlushM}{1} & \sigk{TakenE}{0} & \sig{FlushLane1E} \\
\bottomrule
\end{tabular}
\end{table}

Un salto bien predicho puede coexistir con un stall. Supongamos que el par
(\ins{lw}, \ins{beq}) se emite junto y el \ins{beq} se predice tomado, de modo que Decode
ya contiene el par del destino, cuya primera instrucción usa el registro que carga el
\ins{lw}. Como no hay error de predicción, \sigk{LoadUseD}{0} detiene el par durante un
ciclo, exactamente como en el dual-issue de partida.

\begin{hdl}{Stalls, reagrupamiento y flushes con predicción}
assign lwStallD = LoadUseD_0 & ~MispredictE;          // a misprediction has priority
assign StallD   = lwStallD;
assign StallF   = lwStallD;
assign RegroupD = IssueOneD & ~lwStallD & ~MispredictE;
assign FlushD   = MispredictE;
assign FlushE_0 = lwStallD | MispredictE;
assign FlushE_1 = lwStallD | MispredictE | IssueOneD;
assign FlushM_1 = FlushLane1E;                        // mispredicted in lane 0
\end{hdl}

% =====================================================================
\section{Ejecución con predicción}
\label{sec:exec}

\paragraph{Salto en el slot 0.} La Fig.~\ref{fig:pd-slot0} muestra el bucle de la
Tabla~\ref{tab:loop} a partir de la segunda iteración. En la primera, el \ins{beq} no está
en la tabla, se predice no tomado y se comporta igual que en la Fig.~\ref{fig:pd-nopred};
al resolverse, crea su entrada con el contador en \texttt{10}. Desde entonces:
\begin{itemize}
  \item \textbf{Ciclo~3.} El par (32,~36) se busca y el BTB predice el \ins{beq} como
        tomado (\sigk{PredF}{0}). El siguiente PC es $\sigk{TargetF}{0} = 16$, y el
        \ins{addi x5} del slot~1 se anula (\sigk{AnnulF}{1}).
  \item \textbf{Ciclo~4.} Fetch ya busca la iteración siguiente.
  \item \textbf{Ciclo~5.} La Branch Unit confirma que el salto se toma y que el destino
        es~16. No hay flush, y el contador pasa a \texttt{11}.
\end{itemize}
La iteración baja de cinco a tres ciclos. El único costo que queda es la posición del
slot~1 anulado.

\begin{figure}[htbp]
\centering
\includegraphics[width=0.92\linewidth]{pd-pred-slot0}
\caption{El mismo bucle con predicción: el \ins{beq} del slot~0 se predice tomado en Fetch,
el slot~1 se anula y la iteración dura tres ciclos}
\label{fig:pd-slot0}
\end{figure}

\paragraph{Salto en el slot 1.} Reordenemos el bucle en cuatro instrucciones, de modo que
el \ins{beq} quede en el slot~1 (la inicialización es la misma):
\begin{center}\small
\begin{tabular}{@{}rll@{}}
16 & \ins{loop: slt x4, x1, x2} & \# par (16, 20) \\
20 & \ins{addi x1, x1, 1} & \\
24 & \ins{addi x10, x10, 3} & \# par (24, 28) \\
28 & \ins{beq  x4, x3, loop} & \# salto en el slot 1 \\
\end{tabular}
\end{center}
Ahora el \ins{slt} compara el valor de \ins{x1} anterior al incremento, por lo que ninguna
instrucción depende de su compañera de par y todos los pares se emiten completos. (Por la
misma razón el bucle da una iteración más, once en vez de diez; para la explicación eso no
importa.) La
Fig.~\ref{fig:pd-slot1} muestra el bucle en régimen. Cuando se busca el par (24,~28), el
BTB predice el \ins{beq} del slot~1 (\sigk{PredF}{1}) y el siguiente PC es
$\sigk{TargetF}{1} = 16$. Esta vez no se anula nada: el slot~0 (\ins{addi x10}) es más
antiguo que el salto y se ejecuta siempre. Cada iteración ocupa dos ciclos y el bucle avanza a
dos instrucciones por ciclo, el máximo del procesador. Es además el bucle de dos pares de la
sección~\ref{sec:btb}: en el ciclo~4, la segunda instancia del \ins{beq} se busca mientras
la primera está en Execute.

\begin{figure}[htbp]
\centering
\includegraphics[width=0.80\linewidth]{pd-pred-slot1}
\caption{Salto predicho tomado en el slot~1: el par completo es útil y el bucle avanza a
dos instrucciones por ciclo}
\label{fig:pd-slot1}
\end{figure}

\paragraph{Error de predicción.} En la última iteración, el contador está en \texttt{11}
y el \ins{beq} vuelve a predecirse tomado, pero esta vez no se toma
(Fig.~\ref{fig:pd-mispredict}). En el ciclo~5 la Branch Unit activa \sig{MispredictE}:
el par de Decode y el par de Fetch, que eran la iteración predicha, se descartan, y Fetch
retoma en $\sig{PCCorrectE} = \sig{PCPlus4E} = 36$, es decir, en el \ins{addi x5} que se
había anulado. El contador baja a \texttt{10}, de modo que si el bucle se vuelve a ejecutar,
su primer salto se predice bien. Es la misma penalización de dos ciclos que pagaba el
dual-issue, pero ahora solo una vez por bucle, al salir, y no en cada iteración (más el
fallo de la primera ejecución, cuando el salto todavía no está en la tabla).

\begin{figure}[htbp]
\centering
\includegraphics[width=0.92\linewidth]{pd-mispredict}
\caption{Error de predicción a la salida del bucle: se descartan dos pares y Fetch retoma
en el $\sig{PC}+4$ del \ins{beq}}
\label{fig:pd-mispredict}
\end{figure}

\begin{landscape}
\begin{figure}[p]
\centering
\includegraphics[width=\linewidth,height=0.88\textheight,keepaspectratio]{proc-btb}
\caption{Procesador dual-issue con predicción de saltos. En café: BTB, cadena del siguiente
PC, reagrupamiento y anulación de la predicción, \{Pred, Target\} hasta la Branch Unit,
\sig{PCCorrectE} y actualización del BTB}
\label{fig:proc-full}
\end{figure}
\end{landscape}

% =====================================================================
\section{Resumen}
\label{sec:summary}

La Tabla~\ref{tab:summary} resume qué se agregó al procesador dual-issue y qué se
modificó.

\begin{table}[htbp]
\centering\small
\caption{Del dual-issue al dual-issue con predicción}
\label{tab:summary}
\begin{tabular}{@{}l>{\raggedright\arraybackslash}p{11cm}@{}}
\toprule
& Elementos \\
\midrule
Nuevo & BTB (16 entradas, 2 lecturas, 1 escritura), dos multiplexores más en la cadena del
        siguiente PC, multiplexores de reagrupamiento de \{Pred, Target\}, \sigk{AnnulF}{1}
        y la anulación del slot~1 \\
Ampliado & registros IF/ID (66 bits más) e ID/EX (33 bits más por vía); Branch Unit: selección de la
           predicción y de \sig{PCPlus4E}, \sig{MispredictE}, \sig{PCCorrectE} y señales de
           actualización \\
Modificado & Hazard Unit: \sig{PCSrcE} $\to$ \sig{MispredictE} y
             \sigk{TakenE}{0} $\to$ \sig{FlushLane1E} \\
Sin cambios & Issue Unit, forwarding, register file, memoria de datos, unidades de
              control \\
\bottomrule
\end{tabular}
\end{table}

La optimización se basa en tres ideas:
\begin{enumerate}
  \item \textbf{Adivinar en Fetch, verificar en Execute.} El salto se sigue resolviendo
        donde se resolvía. Lo que cambia es que Fetch ya no espera ese resultado: elige el
        siguiente PC con la predicción, y la Branch Unit solo interviene si la predicción
        falla.
  \item \textbf{La predicción es parte de la instrucción.} Viaja con ella por los
        registros de pipeline, se reagrupa con ella y se descarta con ella. Así la
        verificación siempre compara la predicción correcta, sin importar los stalls ni
        los reagrupamientos que haya entre Fetch y Execute.
  \item \textbf{El orden del par sigue mandando.} Un salto predicho en el slot~0 anula al
        slot~1, que es más joven; un salto en el slot~1 conserva al slot~0, que es más
        antiguo. Es la misma regla que usaba la Hazard Unit para los saltos tomados,
        aplicada ahora en Fetch.
\end{enumerate}

\begin{thebibliography}{9}
\bibitem{harris} S.~L. Harris y D.~M. Harris, \emph{Digital Design and Computer
  Architecture, RISC-V Edition}. Morgan Kaufmann, 2021.
\end{thebibliography}

\end{document}
